% Part: first-order-logic 
% Chapter: axiomatic-deduction 
% Section: deduction-theorem

% verification of properties of provability needed for maximally 
% consistent sets in the completeness chapter.

\documentclass[../../../include/open-logic-section]{subfiles}

\begin{document}

\iftag{FOL}
      {\olfileid{fol}{axd}{ded}}
      {\olfileid{pl}{axd}{ded}}

\olsection{De deductiestelling: een aanname in een implicatie zetten}

Een !!{derivation} in een axiomasysteem uitschrijven kost veel werk,
en zo'n !!{derivation} kan moeilijk te vinden zijn. Vaak hoeven we niet
echt een lange rij !!{formula}s op te schrijven. Het is makkelijker om
met een paar eenvoudige resultaten te laten zien dat zo'n
!!{derivation} bestaat. Drie daarvan hebben we al. Met
\olref[ptn]{prop:reflexivity} mogen we $\Gamma \Proves !A$ schrijven
zodra we weten dat $!A \in \Gamma$. Volgens
\olref[ptn]{prop:monotonicity} mogen we aan $\Gamma \Proves !A$ een
extra aanname toevoegen; dan krijgen we $\Gamma \cup \{!B\} \Proves
!A$. En \olref[ptn]{prop:transitivity} zegt: als $\Gamma \Proves !A$
en $!A \Proves !B$, dan $\Gamma \Proves !B$. Nu volgt nog een
eenvoudig resultaat. Het is een ``meta''-versie van modus ponens: dezelfde
stap, maar nu voor uitspraken over hele afleidingen.

\begin{prop}
  \ollabel{prop:mp} Als $\Gamma \Proves !A$ en $\Gamma \Proves !A \lif
  !B$, dan $\Gamma \Proves !B$.
\end{prop}

\begin{proof}
  De volgende drie regels laten zien dat $\{!A, !A \lif !B\} \Proves
  !B$:
  \begin{derivation}
    1. & $!A$ & Aanname\\
    2. & $!A \lif !B$ & Aanname\\
    3. & $!B$ & 1, 2, MP
  \end{derivation}
  Met \olref[ptn]{prop:transitivity} krijgen we dan $\Gamma \Proves
  !B$.
\end{proof}

Het belangrijkste hulpmiddel dat we hier gebruiken, is de
deductiestelling:

\begin{thm}[Deductiestelling: zet de extra aanname in de implicatie]
\ollabel{thm:deduction-thm} $\Gamma \cup \{!A\} \Proves !B$ dan en
slechts dan als $\Gamma \Proves !A \lif !B$.
\end{thm}

\begin{proof}
De richting ``als'' is direct. Neem aan dat $\Gamma \Proves !A \lif
!B$. Volgens \olref[ptn]{prop:monotonicity} mogen we de extra aanname
toevoegen en krijgen we $\Gamma \cup \{!A\} \Proves !A \lif !B$.
Volgens \olref[ptn]{prop:reflexivity} geldt ook $\Gamma \cup \{!A\}
\Proves !A$. Met \olref{prop:mp} krijgen we dus $\Gamma \cup \{!A\}
\Proves !B$.

Nu de richting ``slechts als''. We gebruiken inductie naar de lengte
van de !!{derivation} van $!B$ uit $\Gamma \cup \{!A\}$: eerst een
afleiding van lengte één, daarna de stap naar een langere afleiding.

Begin met de inductiebasis: een !!{derivation} van lengte~$1$.
!!^a{derivation} van~$!B$ uit $\Gamma \cup \{!A\}$ met lengte~$1$
bestaat alleen uit $!B$. Als die afleiding klopt, geldt ofwel
$!B$$\in \Gamma \cup \{!A\}$, of is de formule een axioma. Eerst het
geval dat $!B \in \Gamma$ of een axioma is. Dan $\Gamma \Proves !B$.
Uit \olref[prp]{ax:lif1} hebben we ook $\Gamma \Proves !B \lif (!A
\lif !B)$. Met \olref{prop:mp} krijgen we $\Gamma \Proves !A \lif !B$.
Blijft het geval $!B \in \{ !A\}$. Ook dan $\Gamma \Proves !A \lif
!B$, want de laatste !!{sentence}~$!A \lif !B$ is dan dezelfde als
$!A \lif !A$. Die hebben we al !!{derive}d in
\olref[pro]{ex:identity}.

Nu de inductiestap. Kijk eerst naar het geval waarin !!a{derivation}
van~$!B$ uit $\Gamma \cup \{!A\}$ eindigt met een stap~$!B$ die met
modus ponens is verkregen. (Is dat niet zo, dan geldt $!B \in \Gamma$,
$!B \ident !A$, of is $!B$ een axioma. Dan gebruiken we dezelfde
redenering als in de inductiebasis.) Voor de laatste stap met modus
ponens moeten eerder in de !!{derivation} de stappen $!C \lif !B$ en
$!C$ staan, voor een !!{formula}~$!C$. Dus $\Gamma \cup \{!A\}
\Proves !C \lif !B$ en $\Gamma \cup \{!A\} \Proves !C$. De twee
bijbehorende !!{derivation}s zijn korter. De inductiehypothese geldt
daarom voor allebei. Dat geeft:
\begin{align*}
  & \Gamma \Proves !A \lif (!C \lif !B); \\
  & \Gamma \Proves !A \lif !C.
\end{align*}
Verder geldt
\[
\Gamma \Proves (!A \lif (!C \lif !B)) \lif
((!A\lif !C)  \lif (!A \lif !B)),
\]
volgens \olref[prp]{ax:lif2}. Pas \olref{prop:mp} nu twee keer toe. Dan
krijgen we $\Gamma \Proves !A \lif !B$, en dat moesten we bewijzen.
\end{proof}

Nu is ook te zien waarom \olref[prp]{ax:lif1} en
\olref[prp]{ax:lif2} precies zo zijn gekozen: daardoor geldt de
deductiestelling.

Hier volgen nog enkele nuttige feiten over !!{derivability}. Het
bewijs ervan is een oefening.

\begin{prop}
  \ollabel{prop:derivfacts}
\begin{enumerate}
\item $\Proves (!A \lif !B) \lif ((!B \lif !C)
  \lif (!A \lif !C)$; \ollabel{derivfacts:a}
\item Als $\Gamma \cup \{ \lnot !A\}
  \Proves \lnot !B$, dan $\Gamma \cup \{ !B\} \Proves
  !A$ (contrapositie: keer de richting om en ontken beide kanten);
  \ollabel{derivfacts:b}
\item  $\{ !A, \lnot!A\} \Proves
    !B$ (ex falso quodlibet, ofwel explosie: uit een tegenspraak volgt
    elke formule); \ollabel{derivfacts:c}
\item  $\{ \lnot\lnot!A\} \Proves
  !A$ (eliminatie van dubbele negatie: verwijder de twee negaties);
  \ollabel{derivfacts:d}
\item Als $\Gamma \Proves \lnot\lnot!A$, dan $\Gamma \Proves
  !A$;\ollabel{derivfacts:e}
\end{enumerate}
\end{prop}

\tagprob{FOL}
\begin{prob}
  Bewijs \olref[fol][axd][ded]{prop:derivfacts}
\end{prob}
\tagendprob

\tagprob{notFOL}
\begin{prob}
  Bewijs \olref[pl][axd][ded]{prop:derivfacts}
\end{prob}
\tagendprob

\end{document}
